An animal that survives an episode may have learned to forage, started near
abundant food, inherited a useful policy, or simply avoided encounters. A stable
population can similarly reflect effective reproduction or a favorable energy
budget rather than improved decision making. Separating these explanations
requires explicit world rules, independent learner state, controlled comparisons
and a recording format that supports retrospective inspection.

\framework\ addresses these practical requirements in an inspectable ecosystem
simulator. Each living animal owns its controller, random streams, optimizer and
experience. Named species select a resource-eating or prey-eating role and a
brain template through JSON. Animals interact through shared physical state
while learning independently. Their bodies can die; eligible pairs can choose
to reproduce; descendants inherit learned parameters without inheriting their
parents' optimizer or experience buffers. This combination produces a changing,
nonstationary population rather than a fixed set of interchangeable agents.

The framework separates simulation, training, recording and visualization.
Headless Python runs produce portable artifacts; JavaScript replays saved data
without loading live controllers. Engineering choices therefore affect both
experimental cost and what evidence remains available afterward. A compressed
file alone is insufficient if the viewer must decode an entire experiment, or
if model-only exports are mistaken for exact resumable checkpoints.

The contributions described here are: (i) a configurable environment and
controller interface for individual learning with changing populations;
(ii) explicit separation of within-lifetime updates, paired parameter
inheritance, and between-run champion initialization; and (iii) a reproducible
recording/checkpoint architecture with native CPU acceleration and lossless,
deduplicated storage. These are system-design contributions. The present draft
reports historical engineering checks and measurements, while population-level
learning outcomes remain to be established by the studies in
Supplementary Section~\ref{sec:observations}.

Multi-agent platforms such as PettingZoo~\citep{terry2021pettingzoo} emphasize
reusable environment interfaces, while MAgent~\citep{zheng2017magent} supports
large interacting populations. \framework\ instead focuses on transparent
per-animal learner ownership, reproductive inheritance and replay artifacts for
small CPU-oriented experiments. It does not claim comparable scale or direct
compatibility with those APIs, and no comparative performance study against
those platforms is included.

The learner implementations draw on established methods: deep Q-learning
\citep{mnih2013dqn}, policy gradients~\citep{williams1992reinforce},
actor--critic learning~\citep{mnih2016asynchronous}, PPO~\citep{schulman2017ppo},
generalized advantage estimation~\citep{schulman2015gae}, distributional quantile
learning~\citep{dabney2018quantile}, action masking~\citep{huang2020masking}, and
centralized-critic PPO~\citep{yu2022mappo}. Their inclusion is not a claim of a
new reinforcement-learning algorithm. In particular, the GRPO controller is a
simulator adaptation of the group-relative objective introduced by
\citet{shao2024deepseekmath}, with local observations and matched hypothetical
world trajectories rather than language-model completions.
